\documentclass[aspectratio=169]{beamer}
\usetheme[numbering=fraction, progressbar=frametitle, block=fill]{metropolis}
\usepackage[T2A]{fontenc}
\usepackage[utf8]{inputenc}
\usepackage[russian]{babel}
\usepackage{graphicx}
\usepackage{booktabs}
\usepackage{listings}
\usepackage{tikz}
\usepackage{xcolor}

\definecolor{accent}{RGB}{52, 136, 255}
\definecolor{green}{RGB}{102, 221, 136}
\definecolor{orange}{RGB}{255, 204, 68}
\definecolor{graybg}{RGB}{30, 30, 50}
\definecolor{red}{RGB}{255, 102, 102}

\setbeamercolor{frametitle}{bg=accent!20, fg=accent}
\setbeamercolor{block title}{bg=accent!20, fg=accent}
\lstset{basicstyle=\ttfamily\footnotesize, backgroundcolor=\color{graybg!10}, breaklines=true}

\title{Antigravity DEX}
\subtitle{Децентрализованная биржа деривативов --- техническая архитектура}
\date{2026}

\begin{document}

\begin{frame}
\maketitle
\centering{\small PBFT консенсус  \quad LMAX Disruptor \quad Arbitrum + TON бриджи}
\end{frame}

\begin{frame}{Проблема}
\begin{itemize}
    \item \textbf{Centralized risk:} регулятор может заблокировать любого CEX оператора
    \item \textbf{Single point of failure:} у HyperLiquid единый CloudFront Gateway ---~2 даунтайма в 2025
    \item \textbf{Закрытый код:} нельзя проверить честность matching-а
    \item \textbf{Стоимость:} ~\$360k/мес на инфраструктуру HyperLiquid
\end{itemize}
\end{frame}

\begin{frame}{Решение: децентрализация без единой точки отказа}

\begin{columns}[t]
\begin{column}{0.48\textwidth}
\textbf{Архитектура}
\begin{itemize}
    \item 3 валидатора, PBFT (2/3 голосов)
    \item Каждый валидатор = свой REST API
    \item Клиент выбирает любой, failover при падении
    \item Open-source код
\end{itemize}
\end{column}
\begin{column}{0.48\textwidth}
\textbf{Преимущества}
\begin{itemize}
    \item Нет единой точки отказа
    \item \$600/мес vs \$360k/мес (HyperLiquid)
    \item Устойчивость к регуляторам
    \item(\dq протокол, не биржа\dq)
\end{itemize}
\end{column}
\end{columns}
\end{frame}

\begin{frame}{Трёхслойная архитектура}

\centering
\begin{tikzpicture}[node distance=1.2em, auto]
\tikzstyle{layer}=[rectangle, rounded corners, minimum width=22em, minimum height=2em, font=\bfseries]
\tikzstyle{comp}=[rectangle, rounded corners, draw, minimum width=7em, minimum height=1.6em, font=\small]

\node[layer, fill=blue!30] (l1) {L1: Ethereum};
\node[layer, fill=blue!20, below=of l1] (l2) {L2: Arbitrum (rollup batches + USDC)};
\node[draw=none, right=3em of l2, font=\small] (tonnote) {\textcolor{gray}{TON: только платежи (Telegram)}};

\node[layer, fill=blue!40, below=of l2] (l3) {L3: Our SMR (3 валидатора)};

\node[comp, below=1.8em of l3] (pbft) {PBFT Consensus};
\node[comp, right=of pbft] (disr) {Disruptor};
\node[comp, right=of disr] (exe) {Execution};
\node[comp, below=of pbft] (wal) {FlatFileStore};

\draw[->, thick] (l1) -- (l2);
\draw[->, thick] (l2) -- (l3);
\draw[->, dashed, gray] (l3.east) -- ++(1.2em,0) |- (tonnote.south);
\draw[->] (pbft) -- (disr);
\draw[->] (disr) -- (exe);
\draw[->] (exe) -- (wal);

\node[font=\tiny, below=0.2em of wal] (note) {Java пакеты l1/ (legacy), l2/ (PBFT), l3/ (persistence) --- не совпадают со слоями};
\end{tikzpicture}
\end{frame}

\begin{frame}{PBFT Consensus}

\begin{columns}[t]
\begin{column}{0.55\textwidth}
\textbf{Раунд консенсуса (1s tick)}
\begin{enumerate}
    \item Leader забирает tx из mempool
    \item PrePrepare ---~рассылает блок
    \item Prepare ---~2/3 валидаторов подтверждают
    \item Commit ---~2/3 финализируют
    \item Блок исполняется локально каждым валидатором
\end{enumerate}
\end{column}
\begin{column}{0.4\textwidth}
\textbf{Характеристики}
\begin{itemize}
    \item Stake-weighted round-robin leader
    \item 1 Byzantine fault (при 3 валидаторах)
    \item Детерминированное исполнение
    \item gRPC (planned) / HTTP (current)
\end{itemize}
\end{column}
\end{columns}
\end{frame}

\begin{frame}{Disruptor Pipeline}

\begin{columns}[t]
\begin{column}{0.48\textwidth}
\textbf{Единственный путь исполнения}
\vspace{0.3em}
\nopagebreak
\begin{itemize}
    \item LMAX Disruptor RingBuffer
    \item Single-thread consumer
    \item No I/O на горячем пути
    \item StateExecutionHandler.onEvent()
    \item Типы: DEPOSIT, WITHDRAW,\newline PLACE\_ORDER, CANCEL\_ORDER,\newline UPDATE\_ORACLE
\end{itemize}
\end{column}
\begin{column}{0.48\textwidth}
\textbf{Производительность}
\vspace{0.3em}
\nopagebreak
\begin{itemize}
    \item[\textcolor{green}{\checkmark}] \textbf{Execution:} $\sim$5M tx/s (один валидатор, in-memory)
    \item[\textcolor{orange}{?}] \textbf{Consensus (3 ноды):} TBD --- нужен network benchmark
    \item[\textcolor{orange}{?}] \textbf{End-to-end:} TBD --- оценка 50--200ms
\end{itemize}
\vspace{1em}
\hfill{\small\color{gray} Benchmark: ThroughputBenchmarkTest}
\end{column}
\end{columns}
\end{frame}

\begin{frame}{OrderBook + Margin + Liquidation}

\begin{columns}[t]
\begin{column}{0.48\textwidth}
\textbf{OrderBook}
\begin{itemize}
    \item Price-time priority
    \item Limit / Market orders
    \item Partial fill support
\end{itemize}
\textbf{MarginManager}
\begin{itemize}
    \item Cross + Isolated margin
    \item Oracle-weighted equity (WIP)
    \item Fast withdrawal (0.3\% fee)
\end{itemize}
\end{column}
\begin{column}{0.48\textwidth}
\textbf{LiquidationEngine}
\begin{itemize}
    \item Isolated: ликвидация по цене
    \item Cross: принудительное закрытие
    \item FundingCalculator: premium index
\end{itemize}
\textbf{Стейты}
\begin{itemize}
    \item Все в памяти (Disruptor)
    \item FlatFileStore: WAL (JSON Lines) + snapshot каждые 30s
    \item Recovery: snapshot $\to$ replay WAL
\end{itemize}
\end{column}
\end{columns}
\end{frame}

\begin{frame}{Мосты: Arbitrum + TON}

\begin{columns}[t]
\begin{column}{0.48\textwidth}
\textbf{Arbitrum = L2 Settlement}
\begin{itemize}
    \item \textbf{Rollup batch posting} (state root + trades)
    \item Депозит: retryable ticket (10--15 мин)
    \item Вывод: 7 дней (standard) / fast 0.3\%
    \item Challenge window: 7 дней
\end{itemize}
\end{column}
\begin{column}{0.48\textwidth}
\textbf{TON = Payment Rail}
\begin{itemize}
    \item \textcolor{gray}{Только депозиты/выводы, нет батчей}
    \item Депозит: $\sim$3 сек (USDT transfer)
    \item Вывод: $\sim$1 сек (immediate)
    \item Telegram Mini App auth
    \item Пользователи tg-* bypass RSA
\end{itemize}
\end{column}
\end{columns}
\vspace{0.5em}
\begin{block}{\textcolor{orange}{Внимание}}
Оба моста --- in-memory mock. Реальная интеграция с Arbitrum Nitro и TON RPC --- следующий шаг.
\end{block}
\end{frame}

\begin{frame}{Persistence: FlatFileStore}

\begin{columns}[t]
\begin{column}{0.48\textwidth}
\textbf{Структура}
\begin{itemize}
    \item WAL: JSON Lines (все ChainTransaction)
    \item Snapshot: полный срез стейта (JSON)
    \item Период: 30s или 10 блоков
    \item Путь: data-\{nodeId\}/l3/
\end{itemize}
\end{column}
\begin{column}{0.48\textwidth}
\textbf{Snapshots содержат}
\begin{itemize}
    \item Balance (free + locked margin)
    \item Open orders
    \item Positions
    \item processedBridgeTicketIds
    \item pendingWithdrawals (isFastWithdraw, feeBps, beneficiary)
\end{itemize}
\end{column}
\end{columns}
\end{frame}

\begin{frame}{Telegram Mini App}

\begin{columns}[t]
\begin{column}{0.48\textwidth}
\textbf{Функционал}
\begin{itemize}
    \item Компактный UI внутри Telegram
    \item WebApp SDK auth (HMAC initData)
    \item Депозит USDT (TON)
    \item Размещение ордеров
    \item Вывод USDT (TON)
    \item Graceful fallback вне Telegram
\end{itemize}
\end{column}
\begin{column}{0.48\textwidth}
\textbf{Файлы}
\begin{itemize}
    \item ui/telegram.html
    \item ui/telegram-app.js
    \item ui/style.css
    \item ui/index.html (полный dashboard)
\end{itemize}
\textbf{Флоу}
\begin{enumerate}
    \item User отправляет USDT в vault
    \item TonBridgePoller $\to$ Mempool $\to$ PBFT
    \item StateExecutionHandler.credit()
    \item Trade $\to$ Withdraw $\to$ TonBridge.sendToWallet()
\end{enumerate}
\end{column}
\end{columns}
\end{frame}

\begin{frame}{Статус проекта}

\begin{columns}[t]
\begin{column}{0.48\textwidth}
\textcolor{green}{\textbf{Сделано}}
\begin{itemize}
    \item PBFT consensus (3 валидатора, 2/3)
    \item Disruptor pipeline
    \item OrderBook (price-time)
    \item MarginManager (cross + isolated)
    \item Liquidation + Funding
    \item ArbitrumBridge mock
    \item TonBridge mock
    \item FlatFileStore (WAL + snapshot)
    \item Telegram Mini App
    \item 36 тестов (17--20s)
\end{itemize}
\end{column}
\begin{column}{0.48\textwidth}
\textcolor{orange}{\textbf{Открытые вопросы (UNRESOLVED.md)}}
\begin{enumerate}
    \item Currency model (single USD / two-pool / oracle-weighted)
    \item Real Arbitrum Nitro + TON RPC
    \item TON Vault custody (multisig / contract)
    \item Validator stake \& slashing
    \item Cross-token withdrawal policy
\end{enumerate}
\vspace{1em}
\textcolor{orange}{\textbf{Что нужно сейчас}}
\begin{itemize}
    \item Security audit
    \item Real bridge integration
    \item Validator economics design
\end{itemize}
\end{column}
\end{columns}
\end{frame}

\begin{frame}{Код}

\textbf{Пакеты} src/main/java/com/example/dex/

\centering
\begin{tabular}{ll}
\toprule
\textbf{Пакет} & \textbf{Назначение} \\
\midrule
server/ & ValidatorNode (PBFT entry point) \\
disruptor/ & LMAX Disruptor pipeline \\
matching/ & OrderBook (price-time) \\
margin/ & Margin, Liquidation, Funding \\
bridge/ & Arbitrum + TON bridges \\
l2/ & PBFT consensus, Mempool, Network \\
l3/persistence/ & FlatFileStore \\
models/ & Order, Trade, ChainTransaction \\
client/ & DexClient SDK \\
cryptography/ & RSA signing \\
\bottomrule
\end{tabular}

\vspace{0.5em}
\centering{\small \textcolor{gray}{Всего: 36 тестов, все проходят. \quad Java 25. \quad Maven single module.}}
\end{frame}

\begin{frame}{Как начать}

\begin{enumerate}
    \item \texttt{git clone https://github.com/rgobov/dex-on-java}
    \item \texttt{mvn compile}
    \item \texttt{mvn test} \quad (36 тестов, $\sim$20s)
    \item Три ноды:
    \begin{itemize}
        \item \texttt{PEERS=val-1:8001,val-2:8002,val-3:8003 mvn exec:java -Dexec.args="val-1 8001 data-val1"}
        \item \texttt{PEERS=... mvn exec:java -Dexec.args="val-2 8002 data-val2"}
        \item \texttt{PEERS=... mvn exec:java -Dexec.args="val-3 8003 data-val3"}
    \end{itemize}
    \item Telegram Mini App: \texttt{http://localhost:8001/telegram.html}
\end{enumerate}
\end{frame}

\end{document}
